\IfFileExists{stacks-project.cls}{%
\documentclass{stacks-project}
}{%
\documentclass{amsart}
}
\usepackage{amssymb}

% For dealing with references we use the comment environment
\usepackage{verbatim}
\usepackage{fancyvrb}
\newenvironment{reference}{\comment}{\endcomment}
%\newenvironment{reference}{}{}
\newenvironment{slogan}{\comment}{\endcomment}
\newenvironment{history}{\comment}{\endcomment}

% For source figures retained as assets
\usepackage{graphicx}

% For commutative diagrams we use Xy-pic
\usepackage[all]{xy}

% We use 2cell for 2-commutative diagrams.
\xyoption{2cell}
\UseAllTwocells

% We use multicol for the list of chapters between chapters
\usepackage{multicol}

% This is generally recommended for better output
\usepackage{lmodern}
\usepackage[T1]{fontenc}
\usepackage{textalpha}

% For cross-file-references
% We prefer the snapshot tmp/xref/#2.aux,
% which latexmkrc replaces atomically, so
% that documents compiled simultaneously
% never read a half-written .aux file.
% The last argument keeps links to #2.pdf.
\usepackage{xr-hyper}
\newcommand{\maybeexternaldocument}[2]{%
  \IfFileExists{#2.aux}{%
    \externaldocument[#1]{#2}%
  }{%
    \PackageWarning{stack}{Missing #2.aux; external references with prefix #1 unavailable}%
  }%
}
\let\externaldocumentorig\externaldocument
\renewcommand{\externaldocument}[2][]{%
  \IfFileExists{tmp/xref/#2.aux}{%
    \externaldocumentorig[#1]{tmp/xref/#2}[#2.pdf]%
  }{%
  \IfFileExists{#2.aux}{%
    \externaldocumentorig[#1]{#2}%
  }{%
    \PackageWarning{stack}{Missing #2.aux; external references with prefix #1 unavailable}%
  }%
  }%
}

% Package for hypertext links:
\usepackage{hyperref}

% For any local file, say "hello.tex" you want to link to please
% use \externaldocument[hello-]{hello}
\externaldocument[introduction-]{introduction}
\externaldocument[conventions-]{conventions}
\externaldocument[sets-]{sets}
\externaldocument[categories-]{categories}
\externaldocument[topology-]{topology}
\externaldocument[sheaves-]{sheaves}
\externaldocument[sites-]{sites}
\externaldocument[stacks-]{stacks}
\externaldocument[fields-]{fields}
\externaldocument[algebra-]{algebra}
\externaldocument[brauer-]{brauer}
\externaldocument[homology-]{homology}
\externaldocument[derived-]{derived}
\externaldocument[simplicial-]{simplicial}
\externaldocument[more-algebra-]{more-algebra}
\externaldocument[smoothing-]{smoothing}
\externaldocument[modules-]{modules}
\externaldocument[sites-modules-]{sites-modules}
\externaldocument[injectives-]{injectives}
\externaldocument[cohomology-]{cohomology}
\externaldocument[sites-cohomology-]{sites-cohomology}
\externaldocument[dga-]{dga}
\externaldocument[dpa-]{dpa}
\externaldocument[sdga-]{sdga}
\externaldocument[hypercovering-]{hypercovering}
\externaldocument[schemes-]{schemes}
\externaldocument[constructions-]{constructions}
\externaldocument[properties-]{properties}
\externaldocument[morphisms-]{morphisms}
\externaldocument[coherent-]{coherent}
\externaldocument[divisors-]{divisors}
\externaldocument[limits-]{limits}
\externaldocument[varieties-]{varieties}
\externaldocument[topologies-]{topologies}
\externaldocument[descent-]{descent}
\externaldocument[perfect-]{perfect}
\externaldocument[more-morphisms-]{more-morphisms}
\externaldocument[flat-]{flat}
\externaldocument[groupoids-]{groupoids}
\externaldocument[more-groupoids-]{more-groupoids}
\externaldocument[etale-]{etale}
\externaldocument[chow-]{chow}
\externaldocument[intersection-]{intersection}
\externaldocument[pic-]{pic}
\externaldocument[weil-]{weil}
\externaldocument[adequate-]{adequate}
\externaldocument[dualizing-]{dualizing}
\externaldocument[duality-]{duality}
\externaldocument[discriminant-]{discriminant}
\externaldocument[derham-]{derham}
\externaldocument[local-cohomology-]{local-cohomology}
\externaldocument[algebraization-]{algebraization}
\externaldocument[curves-]{curves}
\externaldocument[resolve-]{resolve}
\externaldocument[models-]{models}
\externaldocument[functors-]{functors}
\externaldocument[equiv-]{equiv}
\externaldocument[pione-]{pione}
\externaldocument[etale-cohomology-]{etale-cohomology}
\externaldocument[proetale-]{proetale}
\externaldocument[relative-cycles-]{relative-cycles}
\externaldocument[more-etale-]{more-etale}
\externaldocument[trace-]{trace}
\externaldocument[crystalline-]{crystalline}
\externaldocument[spaces-]{spaces}
\externaldocument[spaces-properties-]{spaces-properties}
\externaldocument[spaces-morphisms-]{spaces-morphisms}
\externaldocument[decent-spaces-]{decent-spaces}
\externaldocument[spaces-cohomology-]{spaces-cohomology}
\externaldocument[spaces-limits-]{spaces-limits}
\externaldocument[spaces-divisors-]{spaces-divisors}
\externaldocument[spaces-over-fields-]{spaces-over-fields}
\externaldocument[spaces-topologies-]{spaces-topologies}
\externaldocument[spaces-descent-]{spaces-descent}
\externaldocument[spaces-perfect-]{spaces-perfect}
\externaldocument[spaces-more-morphisms-]{spaces-more-morphisms}
\externaldocument[spaces-flat-]{spaces-flat}
\externaldocument[spaces-groupoids-]{spaces-groupoids}
\externaldocument[spaces-more-groupoids-]{spaces-more-groupoids}
\externaldocument[bootstrap-]{bootstrap}
\externaldocument[spaces-pushouts-]{spaces-pushouts}
\externaldocument[spaces-chow-]{spaces-chow}
\externaldocument[groupoids-quotients-]{groupoids-quotients}
\externaldocument[spaces-more-cohomology-]{spaces-more-cohomology}
\externaldocument[spaces-simplicial-]{spaces-simplicial}
\externaldocument[spaces-duality-]{spaces-duality}
\externaldocument[formal-spaces-]{formal-spaces}
\externaldocument[restricted-]{restricted}
\externaldocument[spaces-resolve-]{spaces-resolve}
\externaldocument[formal-defos-]{formal-defos}
\externaldocument[defos-]{defos}
\externaldocument[cotangent-]{cotangent}
\externaldocument[examples-defos-]{examples-defos}
\externaldocument[algebraic-]{algebraic}
\externaldocument[examples-stacks-]{examples-stacks}
\externaldocument[stacks-sheaves-]{stacks-sheaves}
\externaldocument[criteria-]{criteria}
\externaldocument[artin-]{artin}
\externaldocument[quot-]{quot}
\externaldocument[stacks-properties-]{stacks-properties}
\externaldocument[stacks-morphisms-]{stacks-morphisms}
\externaldocument[stacks-limits-]{stacks-limits}
\externaldocument[stacks-cohomology-]{stacks-cohomology}
\externaldocument[stacks-perfect-]{stacks-perfect}
\externaldocument[stacks-introduction-]{stacks-introduction}
\externaldocument[stacks-more-morphisms-]{stacks-more-morphisms}
\externaldocument[stacks-geometry-]{stacks-geometry}
\externaldocument[moduli-]{moduli}
\externaldocument[moduli-curves-]{moduli-curves}
\externaldocument[examples-]{examples}
\externaldocument[exercises-]{exercises}
\externaldocument[guide-]{guide}
\externaldocument[desirables-]{desirables}
\externaldocument[coding-]{coding}
\externaldocument[obsolete-]{obsolete}
\externaldocument[fdl-]{fdl}
\externaldocument[index-]{index}

% Heap Project documents
\externaldocument[k3-]{k3}
\externaldocument[torelli-]{torelli}
\externaldocument[linear-algebra-]{linear-algebra}
\externaldocument[groups-]{groups}
\externaldocument[complex-analysis-]{complex-analysis}
\externaldocument[functional-analysis-]{functional-analysis}
\externaldocument[categories-]{categories}
\externaldocument[commutative-algebra-]{commutative-algebra}
\externaldocument[ringed-spaces-]{ringed-spaces}
\externaldocument[homological-algebra-]{homological-algebra}
\externaldocument[lie-algebras-]{lie-algebras}
\externaldocument[representation-theory-]{representation-theory}
\externaldocument[smooth-manifolds-]{smooth-manifolds}
\externaldocument[differential-topology-]{differential-topology}
\externaldocument[algebraic-topology-]{algebraic-topology}
\externaldocument[differential-geometry-]{differential-geometry}
\externaldocument[algebraic-geometry-]{algebraic-geometry}
\externaldocument[symplectic-geometry-]{symplectic-geometry}
\externaldocument[complex-geometry-]{complex-geometry}
\externaldocument[hodge-theory-]{hodge-theory}
\externaldocument[%
variations-of-hodge-structures-]{%
variations-of-hodge-structures}
\externaldocument[%
homological-mirror-symmetry-]{%
homological-mirror-symmetry}
\maybeexternaldocument{deformations-}{deformations}
\externaldocument[vertex-algebras-]{vertex-algebras}
\externaldocument[springer-]{springer}
\externaldocument[infty-categories-]{infty-categories}
\externaldocument[seminars-differential-geometry-]{%
seminars-differential-geometry}
\externaldocument[seminars-complex-geometry-]{%
seminars-complex-geometry}
\externaldocument[seminars-algebraic-geometry-]{%
seminars-algebraic-geometry}
\externaldocument[seminars-others-]{%
seminars-others}
\externaldocument[%
geometric-prequantization-]{%
geometric-prequantization}
\externaldocument[surfaces-]{surfaces}
\externaldocument[birational-maps-commutative-algebra-]{birational-maps-commutative-algebra}
\externaldocument[cremona-transformations-]{cremona-transformations}
\externaldocument[derived-categories-of-sheaves-]{derived-categories-of-sheaves}
\externaldocument[geometric-invariant-theory-]{geometric-invariant-theory}
\externaldocument[geometric-stability-conditions-and-group-actions-]{geometric-stability-conditions-and-group-actions}
\externaldocument[moduli-spaces-of-sheaves-]{moduli-spaces-of-sheaves}
\externaldocument[bridgeland-stability-general-theory-]{bridgeland-stability-general-theory}
\externaldocument[hilbert-schemes-]{hilbert-schemes}
\externaldocument[noncommutative-abelian-surfaces-and-kummer-type-hyperkahler-varieties-]{noncommutative-abelian-surfaces-and-kummer-type-hyperkahler-varieties}
\externaldocument[mumford-tate-groups-in-hodge-theory-]{mumford-tate-groups-in-hodge-theory}
\externaldocument[hodge-birational-atoms-2026-]{hodge-birational-atoms-2026}
\externaldocument[xxii-egd-2026-]{%
xxii-egd-2026}
\externaldocument[pde-]{pde}
\externaldocument[probability-]{probability}

\externaldocument[linguistics-]{linguistics}
\externaldocument[genealogia-familiar-]{%
genealogia-familiar}

\externaldocument[%
16th-alga-meeting-2026-]{%
16th-alga-meeting-2026}

\externaldocument[%
iv-jornadas-lefschetz-2026-]{%
iv-jornadas-lefschetz-2026}

\externaldocument[reading-the-masters-]{%
reading-the-masters}

\externaldocument[real-analysis-]{%
real-analysis}

\externaldocument[optimal-transport-]{%
optimal-transport}

\externaldocument[history-]{history}

\externaldocument[seiberg-witten-]{seiberg-witten}

\externaldocument[history-philosophy-biology-]{history-philosophy-biology}

\externaldocument[chemistry-]{chemistry}

\externaldocument[molecular-cell-biology-]{molecular-cell-biology}

\externaldocument[microbiology-]{microbiology}

\externaldocument[protistology-]{protistology}

\externaldocument[mycology-]{mycology}

\externaldocument[botany-]{botany}

\externaldocument[zoology-]{zoology}

\externaldocument[animal-histology-]{animal-histology}

\externaldocument[animal-physiology-]{animal-physiology}

\externaldocument[systematics-biogeography-]{systematics-biogeography}

\externaldocument[ecology-natural-resources-]{ecology-natural-resources}

\externaldocument[earth-sciences-]{earth-sciences}

\externaldocument[biostatistics-]{biostatistics}

\externaldocument[mathematical-systems-biology-]{mathematical-systems-biology}

\externaldocument[cimpa-el-salvador-2024-]{cimpa-el-salvador-2024}

% Theorem environments.
%
\theoremstyle{plain}
\newtheorem{theorem}[subsection]{Theorem}
\newtheorem{proposition}[subsection]{Proposition}
\newtheorem{lemma}[subsection]{Lemma}
\newtheorem{corollary}[subsection]{Corollary}

\theoremstyle{definition}
\newtheorem{definition}[subsection]{Definition}
\newtheorem{example}[subsection]{Example}
\newtheorem{exercise}[subsection]{Exercise}
\newtheorem{situation}[subsection]{Situation}

\theoremstyle{remark}
\newtheorem{remark}[subsection]{Remark}
\newtheorem{remarks}[subsection]{Remarks}

\numberwithin{equation}{subsection}

% Macros
%
\def\lim{\mathop{\mathrm{lim}}\nolimits}
\def\colim{\mathop{\mathrm{colim}}\nolimits}
\def\Spec{\mathop{\mathrm{Spec}}}
\def\Hom{\mathop{\mathrm{Hom}}\nolimits}
\def\Ext{\mathop{\mathrm{Ext}}\nolimits}
\def\SheafHom{\mathop{\mathcal{H}\!\mathit{om}}\nolimits}
\def\SheafExt{\mathop{\mathcal{E}\!\mathit{xt}}\nolimits}
\def\Sch{\mathit{Sch}}
\def\Mor{\mathop{\mathrm{Mor}}\nolimits}
\def\Ob{\mathop{\mathrm{Ob}}\nolimits}
\def\Sh{\mathop{\mathit{Sh}}\nolimits}
\def\NL{\mathop{N\!L}\nolimits}
\def\CH{\mathop{\mathrm{CH}}\nolimits}
\def\proetale{{pro\text{-}\acute{e}tale}}
\def\etale{{\acute{e}tale}}
\def\QCoh{\mathit{QCoh}}
\def\Ker{\mathop{\mathrm{Ker}}}
\def\Im{\mathop{\mathrm{Im}}}
\def\Coker{\mathop{\mathrm{Coker}}}
\def\Coim{\mathop{\mathrm{Coim}}}

% Boxtimes
%
\DeclareMathSymbol{\boxtimes}{\mathbin}{AMSa}{"02}

%
% Macros for moduli stacks/spaces
%
\def\QCohstack{\mathcal{QC}\!\mathit{oh}}
\def\Cohstack{\mathcal{C}\!\mathit{oh}}
\def\Spacesstack{\mathcal{S}\!\mathit{paces}}
\def\Quotfunctor{\mathrm{Quot}}
\def\Hilbfunctor{\mathrm{Hilb}}
\def\Curvesstack{\mathcal{C}\!\mathit{urves}}
\def\Polarizedstack{\mathcal{P}\!\mathit{olarized}}
\def\Complexesstack{\mathcal{C}\!\mathit{omplexes}}
% \Pic is the operator that assigns to X its picard group, usage \Pic(X)
% \Picardstack_{X/B} denotes the Picard stack of X over B
% \Picardfunctor_{X/B} denotes the Picard functor of X over B
\def\Pic{\mathop{\mathrm{Pic}}\nolimits}
\def\Picardstack{\mathcal{P}\!\mathit{ic}}
\def\Picardfunctor{\mathrm{Pic}}
\def\Deformationcategory{\mathcal{D}\!\mathit{ef}}


\begin{document}

\title{Mathematical and Systems Biology}
\maketitle

\phantomsection
\label{section-phantom}

\tableofcontents

\section{What this chapter covers}
\label{section-what-this-chapter-covers}

These notes collect the mathematical side of
the biology degree: population dynamics in
discrete and continuous time, chaos,
competition and metapopulations, a method
for detecting causality in time series, and
the network and gene-circuit models of
systems biology. The mathematics is stated
as mathematics; the biological reading of
each model is given alongside.

\textbf{Sources of these notes.}
Two courses at the Facultad de Ciencias,
UNAM. The first is Biolog\'ia Matem\'atica
(third semester, fall 2014). What survives
is Daniel's commentary on May's paper on the
logistic map, his numerical experiment on
the Lorenz system, and his final paper
``B\'usqueda de causalidad en un modelo de
especies fugitivas'', with its LaTeX sources,
digitized data and Python scripts. The
second is a programming and complex-systems
workshop, filed as Taller (fourth semester,
spring 2015): Python scripts on graphs, a
plant--pollinator network, a final project
computing graph invariants, a short essay on
scales and levels of organization, and a
critique of the ComK--ComS model of
\cite{suel2006excitable}. The notes were
checked and completed with
\cite{murray2002mathematical},
\cite{may1976simple},
\cite{strogatz1994nonlinear}, the original
papers \cite{armstrong1976fugitive} and
\cite{sugihara2012detecting}, the
metapopulation chapter
\cite{nee2007metapopulations}, and the
systems-biology texts
\cite{ideker2001new} and
\cite{stokes2007modeling}.

\section{Exponential and logistic growth}
\label{section-exponential-logistic}

Let $N(t)$ be the size of a population. If
births and deaths are proportional to $N$,
with per capita rates $b$ and $d$, then
$$
\frac{dN}{dt} = (b-d)N,
\qquad
N(t) = N_0 e^{(b-d)t},
$$
so the population grows exponentially if
$b>d$ and dies out if $b<d$ (Malthus).
Verhulst's \emph{logistic} model makes the
per capita growth rate decrease linearly
with density:
\begin{equation}
\label{equation-logistic}
\frac{dN}{dt} = rN\left(1-\frac{N}{K}\right),
\qquad r, K > 0.
\end{equation}
Here $r$ is the intrinsic growth rate and
$K$ the \emph{carrying capacity}. The
equilibria are $N=0$, unstable since the
linearization is $\dot n = rn$, and $N=K$,
stable since the linearization is
$\dot n = -rn$. The solution with
$N(0)=N_0$ is
$$
N(t) =
\frac{N_0 K e^{rt}}{K + N_0(e^{rt}-1)}
\longrightarrow K .
$$
For $N_0<K/2$ the graph is the familiar
sigmoid, with inflection at $N=K/2$. In
general, an equilibrium $N^*$ of
$\dot N = f(N)$ is linearly stable if
$f'(N^*)<0$ and unstable if $f'(N^*)>0$
\cite[Section 1.1]{murray2002mathematical}.

\begin{remark}
\label{remark-logistic-metaphor}
Murray stresses that the logistic equation
is a metaphor for density-dependent
regulation rather than a law: fitting
census data to part of the sigmoid gives
poor predictions
\cite[Section 1.1]{murray2002mathematical}.
\end{remark}

\section{The logistic map and period
doubling}
\label{section-logistic-map}

For populations with non-overlapping
generations (many temperate insects) one
models the size in generation $t+1$ as a
function of the size in generation $t$.
May's canonical example is the
\emph{logistic map}
\begin{equation}
\label{equation-logistic-map}
x_{t+1} = F(x_t) = a\, x_t(1-x_t),
\qquad 0 \le x_t \le 1
\end{equation}
\cite{may1976simple}. Daniel's course notes
on May's paper derive its first properties.

\begin{itemize}
\item $F$ has its maximum $F(1/2)=a/4$, so
$F$ maps $[0,1]$ to itself if and only if
$0\le a\le 4$. For $a>4$ the point $1/2$
leaves $[0,1]$, and an orbit that leaves
$[0,1]$ tends to $-\infty$.
\item The fixed points are the intersections
of the parabola with the diagonal:
$x=0$ and
$$
x^* = 1 - \frac{1}{a}.
$$
For example $a=2.5$ gives $x^*=0.6$, the
value in Daniel's figure.
\item $F'(x^*) = 2-a$. A fixed point is
stable if $|F'(x^*)|<1$, so $x^*$ attracts
the orbits of $(0,1)$ for $1<a<3$; for
$a<1$ every orbit in $[0,1]$ tends to $0$.
\end{itemize}

At $a=3$ the multiplier passes through
$-1$. The second iterate $F^{2}$ then has
two new fixed points,
$$
p, q =
\frac{(a+1)\pm\sqrt{(a+1)(a-3)}}{2a},
$$
which form a 2-cycle of $F$
\cite[Section 2.3]{murray2002mathematical}.
By the chain rule the multiplier of a
$k$-cycle is the product of $F'$ along it.
For the 2-cycle, using $p+q=(a+1)/a$ and
$pq=(a+1)/a^2$,
$$
F'(p)F'(q) = a^2(1-2p)(1-2q)
= -a^2+2a+4,
$$
so the 2-cycle is stable for
$3<a<1+\sqrt{6}\approx 3.449$. At
$1+\sqrt 6$ it gives way to a stable
4-cycle, and so on: a cascade of
period-doubling bifurcations with periods
$2^n$, whose parameter values accumulate at
$a_c \approx 3.5700$ \cite{may1976simple}.

\begin{remark}
\label{remark-feigenbaum}
If $a_{n}$ denotes the parameter of the
bifurcation to period $2^{n}$, then
$$
\lim_{n\to\infty}
\frac{a_{n}-a_{n-1}}{a_{n+1}-a_{n}}
= \delta = 4.6692\ldots,
$$
and Feigenbaum showed that $\delta$ is the
same for all one-humped maps that undergo
period doubling
\cite[Section 2.3]{murray2002mathematical}.
\end{remark}

Beyond $a_c$ the dynamics is
\emph{chaotic} in the sense of Li and Yorke:
there are infinitely many periodic orbits
and uncountably many bounded aperiodic
orbits. The first odd cycle appears at
$a\approx 3.6786$ and the 3-cycle at
$a = 1+\sqrt 8 \approx 3.8284$, after which
there are cycles of every period
\cite{may1976simple}. This is Sharkovskii's
theorem, of which Li and Yorke's ``period
three implies chaos'' is a special case:
for a continuous map of an interval, a
periodic point of period $3$ forces periodic
points of every period
\cite[Section 2.3]{murray2002mathematical}.
Between the chaotic bands there are
windows of stable periodic behaviour, and
the orbit diagram is self-similar.

\begin{remark}
\label{remark-may-message}
May's point for ecologists: apparently
random fluctuations in field data need not
come from environmental noise; a
deterministic density-dependent law with a
steep enough hump produces them
\cite{may1976simple}.
\end{remark}

\section{Lotka--Volterra predator--prey}
\label{section-lotka-volterra}

Let $x$ be the prey and $y$ the predator.
The \emph{Lotka--Volterra} system, in the
notation of Daniel's final paper, is
\begin{equation}
\label{equation-lotka-volterra}
\dot x = ax - bxy,
\qquad
\dot y = -cy + dxy,
\end{equation}
with $a,b,c,d>0$: prey grow exponentially
without predators, predators die out
without prey, and encounters, at rate
proportional to $xy$, transfer biomass.

\textbf{Equilibria and linearization.}
The equilibria are $(0,0)$ and
$(c/d, a/b)$. The Jacobian is
$$
J(x,y) =
\begin{pmatrix}
a-by & -bx \\
dy & -c+dx
\end{pmatrix}.
$$
At the origin the eigenvalues are $a>0$
and $-c<0$: a saddle. At $(c/d,a/b)$,
$$
J =
\begin{pmatrix}
0 & -bc/d \\
ad/b & 0
\end{pmatrix},
\qquad
\lambda = \pm i\sqrt{ac},
$$
a linear centre, with period
$2\pi/\sqrt{ac}$ for small oscillations
\cite[Section 3.1]{murray2002mathematical}.

\textbf{Conserved quantity.}
Linearization alone cannot decide the
nonlinear behaviour at a centre. Here a
first integral does: the function
$$
V(x,y) = dx - c\ln x + by - a\ln y
$$
on the open quadrant satisfies
$$
\dot V
= \left(d-\frac{c}{x}\right)\dot x
+ \left(b-\frac{a}{y}\right)\dot y
= 0 ,
$$
since both terms equal
$\pm(dx-c)(a-by)$. $V$ is strictly convex
with its minimum at $(c/d,a/b)$, so its
level sets are closed curves and every
orbit in the open quadrant is periodic.
Prey and predator oscillate, the predator
peak lagging behind the prey peak.

\begin{example}
\label{example-lv-paper}
Daniel's paper uses $a=c=1$ and $b=d=0.2$.
The equilibrium is $(5,5)$ and small
oscillations have period close to $2\pi$.
The caption of his figure says the
trajectories are ``almost closed'' and
might reach the equilibrium after a long
time; by the conserved quantity they are
exactly closed, and the visible drift is
an artefact of the numerical integrator.
\end{example}

\begin{remark}
\label{remark-lv-structural}
The centre is structurally unstable: any
small change to the vector field (for
instance logistic prey growth) can turn it
into a focus or create a limit cycle. This
is why Lotka--Volterra is a starting point
rather than a realistic model
\cite[Section 3.1]{murray2002mathematical}.
\end{remark}

\section{Competition and coexistence}
\label{section-competition}

Two species with logistic growth inhibit
each other. In the form used in Daniel's
paper,
\begin{equation}
\label{equation-competition}
\begin{aligned}
\dot X_1 &= \frac{r_1}{K_1}\,
X_1\,(K_1 - X_1 - \beta_{12} X_2), \\
\dot X_2 &= \frac{r_2}{K_2}\,
X_2\,(K_2 - \beta_{21} X_1 - X_2),
\end{aligned}
\end{equation}
with all parameters positive;
$\beta_{12}$ measures the effect of species
$2$ on species $1$ and vice versa. This is
Murray's model with $b_{ij}=\beta_{ij}$
\cite[Section 3.5]{murray2002mathematical}.

\textbf{Nullclines.}
Besides the axes, the nullclines are the
lines $X_1 + \beta_{12} X_2 = K_1$ and
$\beta_{21} X_1 + X_2 = K_2$, both with
negative slope. The first meets the axes at
$(K_1,0)$ and $(0,K_1/\beta_{12})$, the
second at $(K_2/\beta_{21},0)$ and
$(0,K_2)$.

\textbf{Equilibria.}
$(0,0)$, $(K_1,0)$, $(0,K_2)$ and, if
$\beta_{12}\beta_{21}\ne 1$,
$$
X^* =
\left(
\frac{K_1-\beta_{12}K_2}
{1-\beta_{12}\beta_{21}},\;
\frac{K_2-\beta_{21}K_1}
{1-\beta_{12}\beta_{21}}
\right).
$$
The origin is an unstable node. At
$(K_1,0)$ the eigenvalues are $-r_1$ and
$r_2(1-\beta_{21}K_1/K_2)$, so species $2$
invades a monoculture of species $1$ if and
only if $K_2 > \beta_{21}K_1$; symmetrically
at $(0,K_2)$.

\textbf{The four cases.}
With $a_{12}=\beta_{12}K_2/K_1$ and
$a_{21}=\beta_{21}K_1/K_2$
\cite[Section 3.5]{murray2002mathematical}:
\begin{enumerate}
\item $a_{12}<1$ and $a_{21}<1$, that is
$K_1/\beta_{12} > K_2$ and
$K_2/\beta_{21} > K_1$: then $X^*$ lies in
the open quadrant and attracts all of it.
This is stable coexistence.
\item $a_{12}>1$ and $a_{21}>1$: $X^*$ is a
saddle, $(K_1,0)$ and $(0,K_2)$ are both
stable, and the separatrix through $X^*$
decides the winner.
\item $a_{12}<1<a_{21}$: species $1$
excludes species $2$.
\item $a_{12}>1>a_{21}$: species $2$
excludes species $1$.
\end{enumerate}
Cases 2--4 are the competitive exclusion
principle. Coexistence requires that each
species limit itself more than it limits
the other.

\begin{example}
\label{example-competition-paper}
Daniel's figure uses $r_1=5$, $r_2=8$,
$K_1=10$, $K_2=15$, $\beta_{12}=0.4$,
$\beta_{21}=0.9$. The conditions hold:
$K_1/\beta_{12}=25>15$ and
$K_2/\beta_{21}\approx 16.7>10$. Solving
$X_1 = 10-0.4X_2$, $X_2 = 15-0.9X_1$ gives
$0.64\,X_1 = 4$, so
$$
X^* = (6.25,\; 9.375),
$$
the intersection shown in his figure. The
Jacobian there is
$$
J =
\begin{pmatrix}
-3.125 & -1.25 \\
-4.5 & -5
\end{pmatrix},
$$
with trace $-8.125$, determinant $10$ and
discriminant
$8.125^2-40\approx 26.0>0$: a stable node.
\end{example}

\begin{remark}
\label{remark-competition-no-cycles}
Planar competition of this kind cannot
oscillate. At $X^*$ the off-diagonal
entries of $J$ are both negative, so the
discriminant $(J_{11}-J_{22})^2 +
4J_{12}J_{21}$ is positive: the eigenvalues
are real and $X^*$ is never a focus.
Globally, with the Dulac function
$B = 1/(X_1X_2)$,
$$
\frac{\partial (B\dot X_1)}{\partial X_1}
+ \frac{\partial (B\dot X_2)}{\partial X_2}
= -\frac{r_1}{K_1X_2}-\frac{r_2}{K_2X_1}
< 0
$$
on the open quadrant, so by the
Bendixson--Dulac criterion there are no
periodic orbits. This is the point behind
Armstrong's remark that oscillations do not
occur in simple two-species competition
models
(Section~\ref{section-metapopulations}).
\end{remark}

\section{The Lorenz system and sensitive
dependence}
\label{section-lorenz}

The \emph{Lorenz system} is
\begin{equation}
\label{equation-lorenz}
\dot x = \sigma(y-x),
\quad
\dot y = \rho x - y - xz,
\quad
\dot z = xy - \beta z,
\end{equation}
with $\sigma,\rho,\beta>0$
\cite[Section 9.2]{strogatz1994nonlinear},
where $\rho$ and $\beta$ are called $r$ and
$b$. The origin is a fixed point, globally
attracting for $\rho<1$. For $\rho>1$ there
are two more,
$$
C^{\pm} =
\left(\pm\sqrt{\beta(\rho-1)},\,
\pm\sqrt{\beta(\rho-1)},\,
\rho-1\right),
$$
linearly stable for
$$
1<\rho<\rho_H =
\frac{\sigma(\sigma+\beta+3)}
{\sigma-\beta-1},
$$
assuming $\sigma>\beta+1$. For
$\sigma=10$, $\beta=8/3$ this gives
$\rho_H = 470/19\approx 24.74$, where $C^\pm$
lose stability in a subcritical Hopf
bifurcation
\cite[Section 9.2]{strogatz1994nonlinear}.

\textbf{Parameter regimes.}
Still with $\sigma=10$, $\beta=8/3$: at
$\rho\approx 13.926$ a homoclinic
bifurcation creates a complicated invariant
set containing infinitely many saddle
cycles; it is not an attractor, but nearby
trajectories wander chaotically for a while
before settling on $C^+$ or $C^-$
(\emph{transient chaos}). At
$\rho\approx 24.06$ this set becomes the
strange attractor; the classical chaotic
value is $\rho=28$
\cite[Section 9.5]{strogatz1994nonlinear}.

\textbf{Daniel's experiment.}
With $\sigma=10$, $\beta=8/3$, $\rho=15$ he
integrated two solutions whose initial $x$
differed by $0.01$. They were
indistinguishable for about 1{,}400
iterations, clearly different by 1{,}800,
and different from then on. With a
difference of $0.1$ the solutions separated
earlier. He concluded that the closer the
initial conditions, the longer they stay
together, and that the system shows
sensitive dependence on initial conditions.
Since $13.926<15<24.74$, this value lies in
the transient-chaos regime: what the
experiment shows is sensitive dependence
near the non-attracting invariant set, and
the two solutions may end on different
equilibria $C^\pm$.

\textbf{Lyapunov exponent and horizon.}
On a chaotic attractor nearby trajectories
separate as
$$
\|\delta(t)\| \sim \|\delta_0\|\,
e^{\lambda t},
$$
with $\lambda\approx 0.9$ for the Lorenz
attractor; $\lambda$ is the largest
Lyapunov exponent. If $a$ is the tolerated
error, prediction fails after
$$
t_{\mathrm{horizon}} \sim
\frac{1}{\lambda}\ln\frac{a}{\|\delta_0\|}
$$
\cite[Section 9.3]{strogatz1994nonlinear}.
So the time to divergence grows only
logarithmically as $\|\delta_0\|$
shrinks: dividing the initial error by $10$
buys a fixed extra time $\ln(10)/\lambda$,
not ten times as much. This refines
Daniel's statement that the period of
similarity is ``proportional'' to the size
of the initial difference.

\section{Metapopulations and fugitive
species}
\label{section-metapopulations}

A \emph{metapopulation} is a population of
local populations on discrete habitat
patches, persisting through colonization of
empty patches and extinction of occupied
ones \cite{nee2007metapopulations}.

\textbf{Levins' model.}
Let $p$ be the fraction of occupied
patches, $c$ the colonization rate and $e$
the local extinction rate. Then
\begin{equation}
\label{equation-levins}
\dot p = c\,p\,(1-p) - e\,p .
\end{equation}
This is logistic growth with
$r = c-e$ and $K = 1-e/c$. If $c>e$ the
equilibrium $p^* = 1-e/c$ is globally
stable on $(0,1]$: a fraction $e/c$ of
perfectly suitable patches is always empty.
If a fraction $1-h$ of the patches is
destroyed, $1$ is replaced by $h$, and the
metapopulation dies out once $h<e/c$
\cite[Section 4.2]{nee2007metapopulations}.

\textbf{Fugitive species.}
A \emph{fugitive species} is an inferior
competitor that persists regionally because
it colonizes newly emptied patches faster
than the dominant species displaces it
(Hutchinson's idea, recalled in
\cite{armstrong1976fugitive}). In the
two-species Levins model where the superior
competitor immediately eliminates the
inferior on arrival, a necessary condition
for coexistence is
$$
\frac{c_i}{e_i} > \frac{c_s}{e_s},
$$
with $s$ and $i$ the superior and inferior
species: the inferior species must have the
larger ``basic reproduction number'' of
patches
\cite[Section 4.4]{nee2007metapopulations}.

\textbf{Armstrong's model.}
Armstrong works with $W$ patches per unit
area, fractions $X_1, X_2$ occupied by the
dominant $S_1$ and the fugitive $S_2$,
effective colonization rates
$c_i = \sigma_i m_i$ (viable propagules per
patch and unit time, times the fraction
landing on patches), a common extinction
rate $\varepsilon$, and a probability
$\phi$ that a propagule of $S_1$ captures a
patch of $S_2$:
$$
\begin{aligned}
\dot X_1 &= X_1\bigl(c_1
[1-X_1-(1-\phi)X_2] - \varepsilon\bigr),\\
\dot X_2 &= X_2\bigl(c_2(1-X_1-X_2)
- \phi c_1 X_1 - \varepsilon\bigr)
\end{aligned}
$$
\cite[equations 1a,
1b]{armstrong1976fugitive}.
This is the competition system
\eqref{equation-competition} with
$$
\begin{gathered}
\frac{r_i}{K_i} = c_i,
\qquad
K_i = 1-\frac{\varepsilon}{c_i},\\
\beta_{12} = 1-\phi,
\qquad
\beta_{21} = 1+\frac{\phi c_1}{c_2}.
\end{gathered}
$$
Stable coexistence holds exactly for
$\varepsilon$ in an interval
$(\varepsilon_{\min},\varepsilon_{\max})$
with
$\varepsilon_{\max}/\varepsilon_{\min}
= c_2/c_1$, which is non-empty if and only
if the better competitor is the poorer
colonizer, $c_1<c_2$. Armstrong calls this
ratio the \emph{coexistence bandwidth}
\cite{armstrong1976fugitive}.

\textbf{The experiment.}
Armstrong's patches were small petri dishes
of agar, colonized by spores blown from
older dishes. \emph{Aspergillus nidulans}
overgrows \emph{Penicillium} colonies (the
better competitor); the \emph{Penicillium}
produces more daughter colonies sooner (the
better colonizer). When all dishes were
renewed at once (ANNUAL regime) the species
never coexisted. When only some dishes were
renewed and new ones were seeded from dishes
of several ages (PERENNIAL regime), they
coexisted, readily at an average dish age
of $2.5$ weeks. Both the percentage area of
\emph{Aspergillus} and the total colony
density then oscillated with a period of
about $10$ weeks, slowly damping.
Armstrong hypothesized that colony density
drives the oscillation (small, dense
colonies are easier for \emph{Aspergillus}
to penetrate) and noted that the system is
not two-dimensional, since dishes of several
ages coexist
\cite{armstrong1976fugitive}. By
Remark~\ref{remark-competition-no-cycles},
no planar competition model of the form
\eqref{equation-competition} reproduces
these oscillations.

\begin{remark}
\label{remark-armstrong-misprint}
Equation (2) of
\cite{armstrong1976fugitive} is printed
without the factors $X_1$ and $X_2$, which
equations (1a), (1b) contain; Daniel's
paper copies the printed form. The factors
do not change the interior equilibrium or
the nullclines, but they are needed for the
axes to be invariant and for the stability
analysis of
Section~\ref{section-competition}.
\end{remark}

\section{Convergent cross mapping}
\label{section-ccm}

\textbf{Granger causality.}
$X$ \emph{Granger-causes} $Y$ if the
predictability of $Y$ drops when $X$ is
removed from the set of possible causal
variables. This presupposes
\emph{separability}: the information about
$X$ is not already contained in the other
variables. Separability holds for purely
stochastic and linear systems but fails in
deterministic dynamics, where, if $X$
drives $Y$, the past of $Y$ already encodes
$X$ \cite{sugihara2012detecting}.

\textbf{Shadow manifolds.}
Given a time series $X(1),\ldots,X(L)$, an
embedding dimension $E$ and a lag $\tau$,
the \emph{shadow manifold} $M_X$ is the set
of delay vectors
$$
x(t) = \bigl(X(t), X(t-\tau), \ldots,
X(t-(E-1)\tau)\bigr).
$$
The justification is Takens' theorem: for
a generic smooth flow on a compact
$d$-dimensional manifold $M$ and a generic
smooth observable, the delay map into
$\mathbb{R}^{2d+1}$ is an embedding (not
checked against Takens' paper). Thus $M_X$
is a diffeomorphic image of the attractor,
and two variables of the same system have
shadow manifolds related by a
diffeomorphism.

\begin{definition}
\label{definition-cross-map}
Let $t_1,\ldots,t_{E+1}$ be the times of
the $E+1$ nearest neighbours of $x(t)$ in
$M_X$, nearest first, and $d_j$ the
Euclidean distance from $x(t)$ to
$x(t_j)$. The \emph{cross-mapped estimate}
of $Y(t)$ from $M_X$ is
$$
\hat Y(t)\mid M_X =
\sum_{j=1}^{E+1} w_j\, Y(t_j),
\qquad
w_j = \frac{u_j}{\sum_k u_k},
$$
where $u_j = \exp(-d_j/d_1)$
\cite{sugihara2012detecting}.
\end{definition}

\textbf{The test.}
If $X$ causally influences $Y$, then $M_Y$
recovers $X$: $\hat X\mid M_Y$ approximates
$X$, and the approximation improves as the
library length $L$ grows and the manifold
fills in. This \emph{convergence}, measured
by the correlation between $\hat X$ and $X$
as a function of $L$, is what distinguishes
causation from correlation. Mutual coupling
gives cross mapping in both directions;
unidirectional forcing $X\Rightarrow Y$
gives convergent estimates of $X$ from
$M_Y$ only, unless the forcing is so strong
that $Y$ synchronizes with $X$. The
standard example: fish abundance can
estimate past climate, but not conversely
\cite{sugihara2012detecting}.

\begin{remark}
\label{remark-ccm-direction}
The direction is easy to invert: ``$M_Y$
estimates $X$'' means ``$X$ causes $Y$''.
In the method section of Daniel's paper the
direction is stated backwards; his
application uses the correct one.
\end{remark}

\section{Daniel's final paper}
\label{section-final-paper}

\textbf{Question.}
Is the total colony density $C$ the cause of
the oscillations in the percentage area $A$
covered by \emph{Aspergillus} in Armstrong's
$2.5$-week PERENNIAL experiment, as
Armstrong suggested? The paper,
``B\'usqueda de causalidad en un modelo de
especies fugitivas'', is a student project
for Biolog\'ia Matem\'atica.

\textbf{Data.}
Armstrong published only figures, so Daniel
digitized them with WebPlotDigitizer. His
scripts hold $41$ weekly values of $A$
(weeks $4$--$44$) and $29$ of $C$; he kept
the $29$ weeks present in both series.

\textbf{Method.}
Shadow manifolds $M_A$ and $M_C$ with
$E=3$ and $\tau=1$ (one sampling step),
plotted in three dimensions with
matplotlib, giving $27$ delay vectors each.
He chose a small cluster of nearby points
in one manifold and marked the points with
the same time indices in the other.

\textbf{Result.}
Four close points in $M_A$ correspond to
close points in $M_C$, so $\hat C\mid M_A$
has small error; a cluster in $M_C$
corresponds to widely scattered points in
$M_A$. By the criterion of
Section~\ref{section-ccm}, $C$ causes $A$
and not conversely, supporting Armstrong's
hypothesis. Following Armstrong, he read
it as density dependence: small, dense
colonies are easy for \emph{Aspergillus}
to penetrate, while large
\emph{Penicillium} colonies resist. He
also suggested that a
density-dependent term coupling both
populations could turn the equilibrium
into a stable focus or a limit cycle.

\begin{remark}
\label{remark-final-paper-caveats}
How far the conclusion goes. (1) With $27$
points no convergence in $L$ can be shown,
and convergence is the defining feature of
the method; the comparison was visual,
without the correlation between $\hat C$
and $C$. (2) After discarding unmatched
weeks the samples are not equally spaced
(for instance weeks $6$, $10$, $11$), so a
lag of one step is not a fixed time lag.
(3) Without convergence, a one-sided
cross-map signal may also come from strong
forcing or a shared driver. The result is
suggestive, not a test. (4) By
Remark~\ref{remark-competition-no-cycles},
the suggested focus or limit cycle cannot
arise in a planar competition model; it
needs either non-competitive terms or more
dimensions, such as the dish ages that
Armstrong pointed to. The paper itself asks
for longer series.
\end{remark}

\section{Networks}
\label{section-networks}

A network is a graph $G=(V,E)$, usually
simple and undirected, with $n=|V|$ and
$m=|E|$. Daniel's workshop scripts compute
the following invariants with the Python
library networkx.

\begin{itemize}
\item The \emph{degree} of $v$ is the size
of its open neighbourhood,
$\deg v = |\{u : uv\in E\}|$; Daniel's
project essay takes it as the basic measure
of a node's role in the system.
\item The \emph{density} is
$2m/(n(n-1))$, the fraction of possible
edges present.
\item The local \emph{clustering
coefficient} of $v$ is
$$
C(v) = \frac{2T(v)}{\deg v(\deg v-1)},
$$
where $T(v)$ is the number of triangles
through $v$: the fraction of pairs of
neighbours of $v$ that are adjacent.
\item \emph{Closeness} of $v$ is
$(n-1)/\sum_{u}d(u,v)$, with $d$ the
graph distance; \emph{betweenness} of $v$
is $\sum_{s\ne v\ne t}
\sigma_{st}(v)/\sigma_{st}$, where
$\sigma_{st}$ counts shortest $s$--$t$
paths and $\sigma_{st}(v)$ those through
$v$.
\end{itemize}

\textbf{Bipartite networks.}
Mutualistic networks such as
plant--pollinator networks are
\emph{bipartite}: $V = P\sqcup A$ with
every edge joining a plant to an animal.
Their \emph{connectance} is $m/(|P||A|)$.
A bipartite graph has no odd cycles, so
every triangle count and ordinary
clustering coefficient is zero;
bipartite-specific versions are needed.

\begin{example}
\label{example-pollination-network}
In the workshop Daniel encoded a
plant--pollinator network of Canary Island
species: plants such as
\emph{Echium plantagineum}, \emph{Laurus
azorica} and \emph{Galactites tomentosa},
and pollinators such as \emph{Bombus
canariensis}, \emph{Apis mellifera} and
\emph{Macroglossum stellatarum}. As encoded
it has $15$ animals, $15$ plants and $74$
interactions, so connectance about $0.33$;
the degrees are very uneven (one unnamed
insect visits $13$ plants, \emph{Galactites}
receives $12$ visitors, several species
have one partner). His script computed
bipartite betweenness and closeness. The
workshop also used a Gal\'apagos pollination
network from the Web of Life database
($18$ species, $25$ interactions,
connectance $0.35$).
\end{example}

\textbf{Final project.}
A program that reads an adjacency matrix
(CSV, from the House of Graphs database),
builds the graph, prints nodes, edges,
degrees, density, triangles and clustering
on request, and draws it.

\textbf{Structure of biological
networks.}
Many biological networks are reported to be
\emph{scale-free}: the probability that a
node has degree $k$ follows a power law,
$\log P(k)$ linear in $\log k$, so a few
hubs carry many links (work of Barab\'asi
and colleagues). Small recurrent subgraphs,
\emph{motifs} such as the feed-forward
loop, and larger \emph{modules} acting as
nearly autonomous functional units are
thought to reflect function. The same
topology can give very different dynamics
depending on parameters, so topology alone
is not enough
\cite{stokes2007modeling}.

\section{An excitable gene circuit}
\label{section-comk-coms}

Under nutrient limitation most cells of
\emph{Bacillus subtilis} sporulate, but a
few percent become \emph{competent} (able
to take up DNA) for some hours and then
return to vegetative growth. The core
circuit
\cite{suel2006excitable}:
\begin{itemize}
\item ComK, the master transcription factor
of competence, activates its own gene
(fast positive feedback);
\item ComK is degraded by the
MecA--ClpP--ClpC protease complex;
\item the peptide ComS competes with ComK
for that protease, so ComS protects ComK;
\item ComK indirectly represses comS
(slow negative feedback through ComS).
\end{itemize}

\textbf{The model.}
In dimensionless form, with $K$ and $S$ the
ComK and ComS levels,
$$
\begin{aligned}
\dot K &= a_k + \frac{b_k K^n}{k_0^n+K^n}
- \frac{K}{1+K+S},\\
\dot S &= \frac{b_s}{1+(K/k_1)^p}
- \frac{S}{1+K+S} + \xi(t),
\end{aligned}
$$
where $\xi$ is noise in ComS expression
\cite[Box 1]{suel2006excitable}. The
shared denominator $1+K+S$ expresses
competition for the protease.

\begin{definition}
\label{definition-excitable}
A dynamical system with a stable
equilibrium is \emph{excitable} if
perturbations below a threshold decay
directly, while perturbations above it
trigger a long excursion through phase
space before the return to the
equilibrium.
\end{definition}

With the parameters of the paper's
supplement ($a_k=0.004$, $b_k=0.07$,
$b_s=0.82$, $k_0=0.2$, $k_1=0.222$, $n=2$,
$p=5$) there are three fixed points: a
stable node at low ComK (vegetative), a
saddle, and an unstable spiral at high
ComK (competent). Recomputing them gives
$K\approx 0.035$, $0.082$, $0.111$. Noise
pushes the state past the saddle's stable
manifold, the trajectory loops around the
spiral (competence) and returns: entry is
random, exit automatic. Experiments
confirmed the predictions: ComK and ComS
promoter activities are anti-correlated
during competence, entry is memoryless
(about $3.6\%$ of division events), and a
rewired strain expressing extra ComS from
a competence promoter enters competence at
the normal rate but mostly cannot exit
\cite{suel2006excitable}.

\begin{remark}
\label{remark-coms-term}
In his workshop critique Daniel questioned
the degradation term $S/(1+K+S)$, which
decreases as $K$ increases and so seems to
contradict the repression of comS by ComK.
There is no contradiction: repression acts
through the production term
$b_s/(1+(K/k_1)^p)$, while the degradation
term says that ComK and ComS compete for
the same protease, so more ComK leaves less
protease for ComS. He reached the right
conclusion that the term does not change
the dynamics: the vegetative state is
stable, and only noise turns the system
on.
\end{remark}

\section{Levels of organization and the
systems approach}
\label{section-systems-approach}

\textbf{Scales and levels.}
In Daniel's workshop essay, a \emph{scale}
classifies measurements by order of
magnitude, while \emph{levels of
organization} (molecule, cell, tissue,
organ, organism, population, community,
ecosystem) form a hierarchy of components
in which each level shows \emph{emergent}
properties absent from the level below.
Both organize phenomena by size or
complexity; the two hierarchies need not
match, since one level can span several
scales.

\textbf{Systems biology.}
Ideker, Galitski and Hood describe systems
biology as studying all elements of a
system and their relations while it
functions, rather than one gene or protein
at a time. Simulations of genetic networks
show emergent properties not visible from
isolated interactions
\cite[p.~351]{ideker2001new}. Their
framework is an iteration
\cite[pp.~354--355]{ideker2001new}:
\begin{enumerate}
\item define all components of the system
and build an initial model;
\item perturb the system systematically
(genetically or environmentally) and
monitor the responses at all levels
(mRNA, protein, pathway);
\item reconcile observed and predicted
responses and refine the model;
\item design new perturbations to
distinguish competing models, and repeat.
\end{enumerate}
The ComK--ComS study
(Section~\ref{section-comk-coms}) follows
this cycle on a small scale: a model of a
core module, single-cell measurements, and
a rewired strain to test a prediction.
Mathematical models of such networks range
from ordinary and partial differential
equations to stochastic and agent-based
models \cite{stokes2007modeling}.

\bibliography{my}
\bibliographystyle{amsalpha}

\end{document}
